\section{Conclusion}

EvoDrive Lab provides a compact benchmark for studying how different learning methods behave on unseen procedural 2D driving tracks. Its main contribution is not realism or scale, but a reproducible and inspectable comparison surface that connects benchmark metrics with visual replay evidence.

Across the completed benchmark, NEAT achieved the strongest held-out completion, the custom genetic algorithm remained highly competitive and benefited from a richer sensor configuration, and the lightweight PPO-style baseline failed to solve the task reliably under the current compute budget. Longer procedural tracks substantially reduced completion for both evolutionary methods, which suggests that track geometry is an important difficulty lever for future versions of the benchmark.

With a fixed experiment protocol, aggregate reporting pipeline, and open-source implementation, EvoDrive Lab can serve as a useful entry point for future work on lightweight driving benchmarks, neuroevolution, and procedural generalization. Natural next steps include stronger policy-gradient baselines, larger procedural suites, and richer qualitative analyses of failure recovery. The main empirical takeaway should still be framed narrowly: these results characterize comparative behavior inside this benchmark rather than broad claims about autonomous driving more generally.
